% !TeX program = xelatex
\documentclass[UTF8,a4paper,zihao=-4,fontset=windows]{ctexart}
\usepackage[margin=21mm,headheight=15pt]{geometry}
\usepackage{amsmath,amssymb,graphicx,xcolor,needspace,longtable,booktabs,fancyhdr}
\usepackage[unicode,hidelinks]{hyperref}
\setlength{\parindent}{0pt}
\setlength{\parskip}{4pt}
\setlength{\emergencystretch}{4em}
\linespread{1.13}\raggedbottom\raggedright
\pagestyle{fancy}\fancyhf{}
\fancyhead[L]{湘教版必修第一册：期中教师资料}
\fancyhead[R]{答案与知识点}\fancyfoot[C]{\thepage}
\graphicspath{{../TeX题库/}}
\newcommand{\subq}[1]{\par\Needspace{2\baselineskip}\noindent（#1）\quad}
\newcommand{\choice}[1]{\par\noindent #1.\quad}
\newcommand{\blank}{\underline{\hspace{1.3cm}}}
\newcommand{\bankfig}[2]{\begin{center}\includegraphics[width=#2\linewidth]{../TeX题库/figures/#1.png}\end{center}}
\newcommand{\itemhead}[2]{\par\Needspace{8\baselineskip}\medskip\hrule\smallskip\textbf{#1}\par{\footnotesize\ttfamily #2}\par}
\newcommand{\field}[2]{\par\textbf{#1：}#2\par}
\hypersetup{pdfauthor={教学资料整理}}

\begin{document}
\begin{center}{\LARGE 期中考试知识点总表}\par\medskip 湘教版必修第一册：至4.2指数函数\end{center}
范围为教材第1—3章、4.1—4.2，至114页。32个知识组关联368个题库条目；以函数为主线，集合与逻辑为工具，不等式为手段。分值30、31、59、30为本次命题约定，120分钟、150分。
本表按所给《教材帮》概念栏目重新归纳，保留其书内页码及PDF位置。章节次序以湘教版为准；L01命题定义另用湘教版14—16页补充。《教材帮》中的对数、导数、三角函数、函数凹凸性等不列作本次新增考查要求。
\tableofcontents\clearpage
\section{集合与逻辑：30分}
\par\medskip\noindent\begin{minipage}{\linewidth}\subsection{S01 集合、元素与常用数集}
集合元素具有确定性、互异性、无序性。元素与集合用 $\in,\notin$；常用数集 $\mathbb N,\mathbb Z,\mathbb Q,\mathbb R$，其中 $0\in\mathbb N$。有限集可含0个元素，空集 $\varnothing$ 不是 $\{0\}$。含参数元素求解后必须复查互异性。
\field{预期方法}{先辨认研究对象和数集范围，再代入定义判断；含参集合逐个来源分类。}
\field{常见错误}{把自然数集误作正整数集；混淆元素和单元素集合；候选参数造成元素重复。}
\field{达标证据}{能解释每个符号的对象类型，并逐个复查候选参数。}
{\small 《教材帮》书内p.1、2；PDF第10、11页。关联13题；例如：教材p.2 1.1 思考 未编号题组；教材p.3 1.1 例题 例题1；教材p.3 1.1 例题 例题2}\par\end{minipage}\par
\par\medskip\noindent\begin{minipage}{\linewidth}\subsection{S02 集合表示、区间与数集点集}
列举法适合元素有限且可逐项确定的集合；描述法写清代表元素及限制。$[a,b]$ 含两端，$(a,b)$ 不含两端，无穷符号一侧只能用圆括号。$\{x\mid\cdots\}$ 的元素是数，$\{(x,y)\mid\cdots\}$ 的元素是有序数对。
\field{预期方法}{先求出约束的解，再选列举法或描述法，必要时在数轴或坐标系中表示。}
\field{常见错误}{遗漏0或负整数；将一个有序数对拆成两个元素；区间端点误取。}
\field{达标证据}{表示方式与原集合完全等价，边界及元素类型均正确。}
{\small 《教材帮》书内p.2、3、4；PDF第11、12、13页。关联12题；例如：教材p.4 1.1 例题 例题3；教材p.5 1.1 例题 例题4；教材p.5 1.1 例题 例题5}\par\end{minipage}\par
\par\medskip\noindent\begin{minipage}{\linewidth}\subsection{S03 子集、真子集与集合相等}
$A\subseteq B$ 表示 $A$ 中每个元素均属于 $B$；真包含还需存在 $B$ 中不属于 $A$ 的元素。$A=B$ 等价于互相包含；空集是任意集合的子集。含 $n$ 个元素的集合有 $2^n$ 个子集、$2^n-1$ 个真子集；非空真子集为 $2^n-2$（$n\ge1$）。
\field{预期方法}{用逐元素检验或数轴包含；参数题先确认是否允许空集，再讨论端点。}
\field{常见错误}{遗漏空集；把元素个数相等当作集合相等；把子集符号当作属于符号。}
\field{达标证据}{能给出包含的双向证据，枚举数量完整。}
{\small 《教材帮》书内p.11、12、13；PDF第20、21、22页。关联11题；例如：教材p.2 1.1 思考 未编号题组；教材p.6 1.1 思考 未编号题组；教材p.7 1.1 例题 例题6}\par\end{minipage}\par
\par\medskip\noindent\begin{minipage}{\linewidth}\subsection{S04 交集、并集及解集的联系}
$x\in A\cap B$ 当且仅当同时属于两集合；$x\in A\cup B$ 当且仅当至少属于一个。交集对应条件“且”，并集对应“或”；乘积为0的解集可取各因式为0的解集之并。数轴上交取公共部分，并取全部覆盖部分。
\field{预期方法}{先各自化简集合或解出约束，再在数轴上合并；有序数对交集通过联立方程求得。}
\field{常见错误}{并集重复写元素；连接处一个闭端点也误删；把交集当并集。}
\field{达标证据}{集合运算与逻辑连接词一致，端点逐个核对。}
{\small 《教材帮》书内p.19、20；PDF第28、29页。关联26题；例如：教材p.9 1.1 思考 未编号题组；教材p.10 1.1 例题 例题9；教材p.10 1.1 例题 例题10}\par\end{minipage}\par
\par\medskip\noindent\begin{minipage}{\linewidth}\subsection{S05 全集、补集与德摩根律}
$\complement_U A=\{x\in U\mid x\notin A\}$，必须先有 $A\subseteq U$。$A\cap\complement_U A=\varnothing$，$A\cup\complement_U A=U$。$\complement_U(A\cup B)=(\complement_U A)\cap(\complement_U B)$；交的补等于补的并。
\field{预期方法}{标出全集内的已知集合，取剩余区域，先算括号内部再补集。}
\field{常见错误}{忘记全集范围；对“且”的否定仍写“且”；漏闭端点。}
\field{达标证据}{每个全集元素恰落在集合及其补集之一。}
{\small 《教材帮》书内p.21；PDF第30页。关联17题；例如：教材p.7 1.1 思考 未编号题组；教材p.8 1.1 例题 例题7；教材p.8 1.1 例题 例题8}\par\end{minipage}\par
\par\medskip\noindent\begin{minipage}{\linewidth}\subsection{S06 Venn图、容斥与实际计数}
有限集合满足 $|A\cup B|=|A|+|B|-|A\cap B|$。三集合容斥先减两两交集再加三重交集。“只参加两项”不同于“两项都参加”；分区人数均须非负。
\field{预期方法}{先画互不重叠区域，从最内层填写人数，再依据总人数或非负限制求解。}
\field{常见错误}{把只参加两项的人数当作完整交集；重复计算三重交集；忽略一种都不参加。}
\field{达标证据}{所有分区非负且按每一种口径加回题设人数。}
{\small 《教材帮》书内p.22；PDF第31页。关联3题；例如：教材p.13 习题1.1 第14题；教材p.13 习题1.1 第17题；教材p.30 复习题一 第15题}\par\end{minipage}\par
\par\medskip\noindent\begin{minipage}{\linewidth}\subsection{L01 命题、真假、逆命题与否定}
命题是可判断真假的陈述句。未限定变量的开放语句不一定是命题。“若 $p$ 则 $q$”的逆命题为“若 $q$ 则 $p$”；否定则是原断言不成立，二者不同。原命题与逆命题的真假必须分别判断。本点的完整定义还依据湘教版教材14—16页；《教材帮》30页主要用于条件命题联系。
\field{预期方法}{明确论域与断言；真命题给理由，假命题给满足题设而使结论失败的反例。}
\field{常见错误}{把逆命题当成否定；只改变结论中的个别词而改变原逻辑范围。}
\field{达标证据}{真假和所给证明或反例严格对应。}
{\small 《教材帮》书内p.30；PDF第39页。关联10题；例如：教材p.14 1.2 思考 未编号题组；教材p.15 1.2 例题 例题1；教材p.15 1.2 例题 例题2}\par\end{minipage}\par
\par\medskip\noindent\begin{minipage}{\linewidth}\subsection{L02 充分、必要、充要条件}
$p\Rightarrow q$ 表示 $p$ 是 $q$ 的充分条件，$q$ 是 $p$ 的必要条件；双向成立为充要。只正向为充分不必要，只反向为必要不充分；两向均不成立为既不充分又不必要。解集关系中 $P\subseteq Q$ 对应 $p\Rightarrow q$；$\neg q\Rightarrow\neg p$ 与之等价。
\field{预期方法}{分别检查两个方向；用解集包含或具体反例判定，不凭句子位置猜方向。}
\field{常见错误}{把“必要”箭头反写；只有一个方向证据便断言“不必要”；忽略给定论域。}
\field{达标证据}{两方向均有证据，严格条件与非严格条件区别明确。}
{\small 《教材帮》书内p.30、31；PDF第39、40页。关联18题；例如：教材p.17 1.2 例题 例题3；教材p.17 1.2 例题 例题4；教材p.18 1.2 例题 例题5}\par\end{minipage}\par
\par\medskip\noindent\begin{minipage}{\linewidth}\subsection{L03 全称、存在量词及真假的判定}
$\forall x\in M,p(x)$ 对所有对象断言；$\exists x\in M,p(x)$ 对至少一个对象断言。全称真要证明任意对象，假只需一个反例；存在真只需一个实例，假要证明所有对象都不满足。多个量词的顺序不能随意交换。
\field{预期方法}{先确认量词、作用对象及论域，再采用任意性证明或实例反例。}
\field{常见错误}{用几个满足的数证明全称命题；把“存在”理解为“唯一存在”。}
\field{达标证据}{证明或例子确在规定论域内。}
{\small 《教材帮》书内p.38、39；PDF第47、48页。关联9题；例如：教材p.19 1.2 思考 未编号题组；教材p.20 1.2 例题 例题6；教材p.20 1.2 思考 未编号题组}\par\end{minipage}\par
\par\medskip\noindent\begin{minipage}{\linewidth}\subsection{L04 含一个量词命题的否定}
$\neg(\forall x\in M,p(x))\Leftrightarrow\exists x\in M,\neg p(x)$；$\neg(\exists x\in M,p(x))\Leftrightarrow\forall x\in M,\neg p(x)$。论域不变，量词交换，断言否定；$\ge$ 的否定为 $<$，$=$ 的否定为 $\ne$。
\field{预期方法}{分三步书写：保留论域、交换量词、否定断言，再判断新命题真假。}
\field{常见错误}{只改量词不改结论；把不等号否定成反向但遗漏等号。}
\field{达标证据}{三个部分全部正确，再给实例或证明核验真假。}
{\small 《教材帮》书内p.39、40；PDF第48、49页。关联10题；例如：教材p.21 1.2 思考 未编号题组；教材p.22 1.2 例题 例题9；教材p.22 1.2 例题 例题10}\par\end{minipage}\par
\section{一元二次函数、方程和不等式：31分}
\par\medskip\noindent\begin{minipage}{\linewidth}\subsection{I01 不等关系与作差比较}
对实数 $A,B$，$A>B\Leftrightarrow A-B>0$。常将差因式分解或配方，利用平方非负定号。差可为0时要保留相等情形；实际不等关系须同时写明字母意义和单位。
\field{预期方法}{作差、化简、定号、回写比较结论。}
\field{常见错误}{看见负平方就写恒小于，遗漏平方为0的情形。}
\field{达标证据}{差的符号及等号条件均与比较结论一致。}
{\small 《教材帮》书内p.44；PDF第53页。关联10题；例如：教材p.32 2.1 思考 未编号题组；教材p.34 2.1 例题 例题1；教材p.34 2.1 例题 例题2}\par\end{minipage}\par
\par\medskip\noindent\begin{minipage}{\linewidth}\subsection{I02 不等式基本性质与符号条件}
不等式同加同减方向不变；同乘除正数保向，负数反向。$a>b,c>d$ 可同向相加；一般不能同向相减相除。正数可保序平方、开方；倒数比较需同号且非0。
\field{预期方法}{每一步明确所乘除因子的正负与非零，或作差求证。}
\field{常见错误}{未知符号就直接乘除；将 $c^2$ 当作必为正数而忘记 $c=0$。}
\field{达标证据}{推理每一步所需条件均由题设保证。}
{\small 《教材帮》书内p.45、46；PDF第54、55页。关联13题；例如：教材p.35 2.1 思考 未编号题组；教材p.37 2.1 例题 例题3；教材p.37 2.1 例题 例题4}\par\end{minipage}\par
\par\medskip\noindent\begin{minipage}{\linewidth}\subsection{I03 基本不等式、等号与最值建模}
$a,b>0$ 时 $a+b\ge2\sqrt{ab}$，等号当且仅当 $a=b$；$u,v\in\mathbb R$ 时 $u^2+v^2\ge2uv$。定积求和最小、定和求积最大；配凑、拆项、换元后必须复查等号是否在定义域内可取。矩形面积、费用等实际题先列约束。
\field{预期方法}{列正数条件与固定量，构造基本不等式，解等号条件并代回实际限制。}
\field{常见错误}{把一个下界误当最小值；多个等号条件不能同时实现；漏单位和整数限制。}
\field{达标证据}{同时写出界、取等点及可取性，实际模型能解释回原情境。}
{\small 《教材帮》书内p.52、53、54；PDF第61、62、63页。关联30题；例如：教材p.33 2.1 案例 案例1；教材p.39 2.1 例题 例题5；教材p.39 2.1 例题 例题6}\par\end{minipage}\par
\par\medskip\noindent\begin{minipage}{\linewidth}\subsection{Q01 二次函数、零点、判别式与根系关系}
$a\ne0$ 时，$ax^2+bx+c=a(x+b/(2a))^2-\Delta/(4a)$，$\Delta=b^2-4ac$。$\Delta>0,=0,<0$ 分别对应两个不同零点、一个零点、无零点。实根满足 $x_1+x_2=-b/a,x_1x_2=c/a$。两正根需根存在、和正、积正；若要求相异还需 $\Delta>0$。
\field{预期方法}{按已知零点、顶点或系数选择交点式、顶点式或一般式；根平方和用根系关系消元。}
\field{常见错误}{遗漏 $a
e0$；把二重根算成两个不同零点；有两正根只查积正。}
\field{达标证据}{参数使所有原条件同时成立，根和积及实际零点个数一致。}
{\small 《教材帮》书内p.65；PDF第74页。关联28题；例如：教材p.16 p.16 练习 第1题；教材p.45 2.2 思考 未编号题组；教材p.46 2.2 例题 例题1}\par\end{minipage}\par
\par\medskip\noindent\begin{minipage}{\linewidth}\subsection{I04 一元二次不等式与参数分类}
$a>0$、两根 $r_1<r_2$ 时，二次式在两根外为正、两根内为负；$a<0$ 反向。严格号排除根，非严格号包含根。$\Delta=0$ 时化平方单独判断，解集可能是空集、单点集、全体实数或挖去一点。含参数时先分二次项为0，再分根的大小。
\field{预期方法}{化标准式、查系数、求根或判别式、按图定号，最后与实际范围取交集。}
\field{常见错误}{套“内小外大”却忘开口；丢重根或降次；漏整数产量限制。}
\field{达标证据}{逐类解集无遗漏，在根及区间内外代入复核。}
{\small 《教材帮》书内p.64、65、66；PDF第73、74、75页。关联35题；例如：教材p.49 习题2.2 第7题；教材p.50 2.3 思考 未编号题组；教材p.52 2.3 例题 例题1}\par\end{minipage}\par
\par\medskip\noindent\begin{minipage}{\linewidth}\subsection{I05 简单分式不等式}
先移项通分，使一边为0。$A/B>0\Leftrightarrow AB>0$（自动排除 $B=0$）；$A/B\ge0\Leftrightarrow AB\ge0$ 且 $B\ne0$。不能不明分母符号就直接两边乘分母；约分前排除的点仍须保留限制。
\field{预期方法}{移项通分、标分子零点和分母禁点、逐段定号。}
\field{常见错误}{非严格不等式将分母零点也收进解集。}
\field{达标证据}{分子零点是否可取、所有分母零点是否排除均写清。}
{\small 《教材帮》书内p.68、69；PDF第77、78页。关联4题；例如：教材p.54 2.3 例题 例题4；教材p.55 p.55 练习 第1题；教材p.59 习题2.3 第5题}\par\end{minipage}\par
\par\medskip\noindent\begin{minipage}{\linewidth}\subsection{I06 恒成立、有解与范围内最值}
$a\ne0$ 时，全域 $ax^2+bx+c>0$ 等价于 $a>0,\Delta<0$；全域 $\ge0$ 则 $a>0,\Delta\le0$。负号对应开口向下。未限定二次时须另查 $a=0$。在区间恒成立转为最小值条件，有解转为相应最大或最小值条件，须检查端点是否可取。
\field{预期方法}{先明确主元及参数，分降次情况，再用判别式、配方或区间最值。}
\field{常见错误}{严格号与非严格号混淆；把有解条件当恒成立；直接用全域判别式处理有限区间。}
\field{达标证据}{所有参数临界值均代回原不等式，不只代回中间式。}
{\small 《教材帮》书内p.71、72、73；PDF第80、81、82页。关联6题；例如：教材p.33 2.1 案例 案例2；教材p.55 2.3 例题 例题5；教材p.55 p.55 练习 第3题}\par\end{minipage}\par
\section{函数概念与性质：59分}
\par\medskip\noindent\begin{minipage}{\linewidth}\subsection{F01 函数本质、构成要素与相同函数}
$A,B$ 为非空实数集，对每个 $x\in A$ 在 $B$ 中恰有一个确定值对应才是函数。允许多对一，不允许一对多或某输入无值；值域可真包含于 $B$。相同函数由定义域和对应关系判断，自变量字母可不同。
\field{预期方法}{逐一核对每个允许输入是否有且仅有一个合乎目标集合的输出。}
\field{常见错误}{把值域直接当作目标集合；只看解析式外观；认为常函数不是函数。}
\field{达标证据}{能指出不成立的具体输入或确认任意输入的唯一输出。}
{\small 《教材帮》书内p.77、78；PDF第86、87页。关联9题；例如：教材p.65 3.1 思考 未编号题组；教材p.65 3.1 思考 未编号题组；教材p.66 3.1 例题 例题2}\par\end{minipage}\par
\par\medskip\noindent\begin{minipage}{\linewidth}\subsection{F02 定义域与原式限制}
分母非0，偶次根式被开方数非负，根式在分母时被开方数严格正；各条件同时满足取交集。零次幂底数非0，负次幂排除0。复合求值还需内层结果进入外层定义域；实际函数加实际范围。
\field{预期方法}{按原式逐项列条件后求交集，不先约分再判断。}
\field{常见错误}{遗漏分母禁点；将函数值范围错当自变量范围；把退化矩形算入实际定义域。}
\field{达标证据}{所有边界、原分母及实际限制逐项代回检查。}
{\small 《教材帮》书内p.79、80、81；PDF第88、89、90页。关联16题；例如：教材p.66 3.1 例题 例题1；教材p.66 3.1 例题 例题2；教材p.67 p.67 练习 第2题}\par\end{minipage}\par
\par\medskip\noindent\begin{minipage}{\linewidth}\subsection{F03 函数求值、解析式与函数建模}
$f(a)$ 是输入 $a$ 的输出，$f(a)+1$ 与 $f(a+1)$ 不同。复合求值由内到外；列表、图象先读取内层值。求解析式可用待定系数、换元、配凑或消元。实际建模顺序为设变量、写范围、列关系、求解、解释。
\field{预期方法}{标清每次输入和输出；按已知函数类型或代换关系建立等价方程。}
\field{常见错误}{将 $f(x)$ 当成 $f	imes x$；复合顺序反；换元后沿用错误定义域。}
\field{达标证据}{表达式满足原函数关系，求值全过程在定义域内，实际单位与范围正确。}
{\small 《教材帮》书内p.80、87、88、89、90、131；PDF第89、96、97、98、99、140页。关联52题；例如：教材p.33 2.1 案例 案例2；教材p.41 2.1 例题 例题8；教材p.41 2.1 例题 例题9}\par\end{minipage}\par
\par\medskip\noindent\begin{minipage}{\linewidth}\subsection{F04 函数三种表示法、图象与读数}
列表法突出数据对应，解析法便于精算，图象法呈现变化趋势。作图先查定义域、截距、顶点或转折点，再连接；空心点不属于图象。读图解 $f(x)>g(x)$ 是取 $f$ 曲线在上方的横坐标。数据间插值只是估计，不唯一确定精确函数。
\field{预期方法}{把关键点与定义域一起标到图上，逐区间读交点和上下关系。}
\field{常见错误}{将示意图预测当事实；把两曲线交点收进严格不等式；折线漏端点。}
\field{达标证据}{关键点、趋势、端点及估计的精度说明都正确。}
{\small 《教材帮》书内p.85、86、93、94；PDF第94、95、102、103页。关联29题；例如：教材p.30 复习题一 第17题；教材p.46 2.2 例题 例题1；教材p.49 习题2.2 第4题}\par\end{minipage}\par
\par\medskip\noindent\begin{minipage}{\linewidth}\subsection{F05 单调性定义、证明与参数}
区间 $I$ 内任意 $x_1<x_2$，若 $f(x_1)<f(x_2)$ 则递增，若相反则递减。证明用取值、作差、变形、定号、结论；两个分离区间同向单调不代表并集上单调。二次函数根据开口与对称轴确定单调区间。
\field{预期方法}{用定义比较任意两值，或依据已知图象性质；含参时使指定区间被相应单调区间覆盖。}
\field{常见错误}{用有限采样代替证明；把分离半轴合并；漏跨零点情形。}
\field{达标证据}{区间、任意性和差的严格符号均写清，参数端点单独核验。}
{\small 《教材帮》书内p.97、98、100；PDF第106、107、109页。关联27题；例如：教材p.72 p.72 练习 第3题；教材p.79 3.2 思考 未编号题组；教材p.81 3.2 思考 未编号题组}\par\end{minipage}\par
\par\medskip\noindent\begin{minipage}{\linewidth}\subsection{F06 最大值、最小值与值域}
最大值需 $f(x)\le M$ 对所有输入成立，且至少一个输入实际取到 $M$；最小值类似。闭区间二次函数比较顶点（在区间内时）与两端点；开端点只有界不一定有最值。值域必须记录所有实际能取到的值。
\field{预期方法}{根据图象、配方、单调性或基本不等式找界，再验证取值点。}
\field{常见错误}{把上界下界当最值；只比较两端漏顶点；只有最值却未说明全部值均可取。}
\field{达标证据}{界和达到该界的点同时提供；区间端点与值域括号一致。}
{\small 《教材帮》书内p.99、83、84；PDF第108、92、93页。关联16题；例如：教材p.72 p.72 练习 第3题；教材p.77 习题3.1 第6题；教材p.79 3.2 思考 未编号题组}\par\end{minipage}\par
\par\medskip\noindent\begin{minipage}{\linewidth}\subsection{F07 分段函数、绝对值及分段应用}
分段函数整体是一个函数。先由输入选择分段再代值；反求输入时须分别求解并与该段范围取交。绝对值按内部为0的点分段。阶梯收费只对超出部分加价；图象各段的端点不能重复给出不同输出。
\field{预期方法}{先列转折点和各段条件，再建立分段式，检查衔接、定义域与值域。}
\field{常见错误}{跨段套公式；超额部分按全量收费；绝对值去号漏区间。}
\field{达标证据}{每段表达式和范围同时正确，临界点的函数值唯一。}
{\small 《教材帮》书内p.86、91、92、93；PDF第95、100、101、102页。关联14题；例如：教材p.74 3.1 例题 例题6；教材p.75 3.1 例题 例题7；教材p.75 3.1 例题 例题8}\par\end{minipage}\par
\par\medskip\noindent\begin{minipage}{\linewidth}\subsection{F08 奇偶性、对称与函数性质应用}
定义域先关于原点对称。$f(-x)=f(x)$ 为偶函数，关于 $y$ 轴对称；$f(-x)=-f(x)$ 为奇函数，关于原点对称。奇函数若0有定义则 $f(0)=0$；零函数可既奇又偶。偶函数对称两侧单调方向相反，奇函数对称两侧相同；这不自动给出全域单调性。
\field{预期方法}{检查定义域，再代入 $-x$；求另一侧解析式时同时写出 $-x$ 所在区间。}
\field{常见错误}{只看公式不看定义域；用一对点证明奇偶；负半轴代换遗漏最外层负号。}
\field{达标证据}{对称范围与恒等式完整，拼接处和跨区间比较均核验。}
{\small 《教材帮》书内p.113、114、115；PDF第122、123、124页。关联29题；例如：教材p.83 3.2 思考 未编号题组；教材p.83 3.2 思考 未编号题组；教材p.84 3.2 例题 例题4}\par\end{minipage}\par
\par\medskip\noindent\begin{minipage}{\linewidth}\subsection{F09 图象平移、伸缩与对称}
向右平移 $h$ 得 $y=f(x-h)$，向上平移 $k$ 得 $y=f(x)+k$；$f(-x)$ 关于 $y$ 轴对称，$-f(x)$ 关于 $x$ 轴对称。$cf(x)$ 将纵坐标乘 $c$，$c<0$ 时同时翻折；水平平移同步改变定义域。
\field{预期方法}{追踪点 $(u,f(u))$ 变换后的坐标，反推新解析式。}
\field{常见错误}{平移只改公式而未改定义域；把纵向倍率当作指数底数变化。}
\field{达标证据}{变换后的关键点、禁点和解析式对应。}
{\small 《教材帮》书内p.86、94；PDF第95、103页。关联6题；例如：教材p.92 复习题三 第19题；教材p.110 4.2 思考 未编号题组；教材p.113 p.113 练习 第1题}\par\end{minipage}\par
\section{幂函数、指数函数：30分}
\par\medskip\noindent\begin{minipage}{\linewidth}\subsection{P01 方根、根式与绝对值}
奇数 $n>1$ 时实数 $a$ 有唯一实 $n$ 次方根；偶数 $n$ 时正数有正负两根，根号表示其中非负的算术根。$(\sqrt[n]a)^n=a$（须有意义）；$\sqrt[n]{a^n}=a$（$n$ 奇），$=|a|$（$n$ 偶）。
\field{预期方法}{先看根指数奇偶，再检查被开方数符号，必要时写绝对值分段。}
\field{常见错误}{把偶次根式值写成正负两个；$\sqrt{a^2}$ 直接写 $a$。}
\field{达标证据}{原根式的非负性和变量条件都得到保留。}
{\small 《教材帮》书内p.143；PDF第152页。关联2题；例如：教材p.95 4.1 例题 例题1；教材p.104 习题4.1 第1题}\par\end{minipage}\par
\par\medskip\noindent\begin{minipage}{\linewidth}\subsection{P02 有理数、实数指数幂与运算}
$a>0$ 时 $a^{m/n}=\sqrt[n]{a^m}$，$a^{-r}=1/a^r$；$a^ra^s=a^{r+s}$，$(a^r)^s=a^{rs}$，$(ab)^r=a^rb^r$（按正底数条件使用）。不能把加法的幂拆成幂的和。无理指数幂按有理指数逼近扩充；根式运算中原分母非0等限制不能约掉。
\field{预期方法}{统一为指数幂，分开处理系数与各底数的指数；最终复查定义条件及数值精度。}
\field{常见错误}{负指数误作负数；幂的幂指数相加；约分丢禁点；过早四舍五入。}
\field{达标证据}{指数代数准确，结果满足原式并保留所有定义限制。}
{\small 《教材帮》书内p.144、145；PDF第153、154页。关联25题；例如：教材p.96 4.1 例题 例题2；教材p.96 4.1 例题 例题3；教材p.96 4.1 例题 例题4}\par\end{minipage}\par
\par\medskip\noindent\begin{minipage}{\linewidth}\subsection{P03 幂函数的概念、图象与性质}
幂函数形如 $y=x^\alpha$，底数是自变量，指数为常数，系数为1。正半轴上 $\alpha>0$ 递增、$\alpha<0$ 递减，均过 $(1,1)$；正指数还可在0取0。全域的定义域与奇偶性须依具体指数判断。重点掌握 $x,x^2,x^3,\sqrt x,1/x$，以及由它们推得的负整数、分数指数实例。
\field{预期方法}{设 $f(x)=x^\alpha$，用已知点求指数，再查该指数的定义域、奇偶性和各单调区间。}
\field{常见错误}{把 $3x^2$ 当标准幂函数；将第一象限性质直接用于负半轴；分数负指数遗漏0禁点。}
\field{达标证据}{表达式、定义域、奇偶性、值域及单调区间彼此一致。}
{\small 《教材帮》书内p.125、126、127；PDF第134、135、136页。关联13题；例如：教材p.102 4.1 例题 例题8；教材p.102 4.1 例题 例题9；教材p.103 p.103 练习 第1题}\par\end{minipage}\par
\par\medskip\noindent\begin{minipage}{\linewidth}\subsection{E01 指数函数的概念、图象与性质}
$y=a^x$ 中 $a>0,a\ne1$，指数是自变量。定义域 $\mathbb R$，值域 $(0,+\infty)$，过 $(0,1)$；$a>1$ 递增，$0<a<1$ 递减。互为倒数的底数对应图象关于 $y$ 轴对称；取 $x=1$ 时纵坐标就是底数。
\field{预期方法}{先查底数条件，利用已知点定底数，再据单调性及图象求值或范围。}
\field{常见错误}{只要求底数非负；漏排1；将 $a^x+b$ 当标准指数函数且仍用原值域。}
\field{达标证据}{底数所有限制、图象定点、单调性与值域完整对应。}
{\small 《教材帮》书内p.150、151；PDF第159、160页。关联20题；例如：教材p.99 4.1 例题 例题7；教材p.100 p.100 练习 第3题；教材p.109 4.2 思考 未编号题组}\par\end{minipage}\par
\par\medskip\noindent\begin{minipage}{\linewidth}\subsection{E02 指数式比较、方程与不等式}
同底比较用指数函数单调性，底数在 $(0,1)$ 时反向；同指数比较可用正半轴幂函数性质；两者都不同时可搭桥或与1比较。含 $a^{|u(x)|}$ 时先分绝对值或由指数非负判断值域。
\field{预期方法}{将比较对象化成同底、同指数，或加入合适中间量；解后复查范围。}
\field{常见错误}{所有底数都按递增；底数和指数同时变化时凭直觉判断。}
\field{达标证据}{每次比较都有明确保序或反序条件，严格性正确。}
{\small 《教材帮》书内p.154、155；PDF第163、164页。关联3题；例如：教材p.111 4.2 例题 例题4；教材p.113 p.113 练习 第3题；教材p.114 习题4.2 第9题}\par\end{minipage}\par
\par\medskip\noindent\begin{minipage}{\linewidth}\subsection{E03 指数增长、衰减与实际模型}
按固定倍率变化的模型 $M(t)=M_0q^t$。增长率 $r$ 对应 $q=1+r$，衰减率 $r$ 对应 $q=1-r$；半衰期 $T$ 可写 $M(t)=M_0(1/2)^{t/T}$。初始时刻通常为0；细胞分裂、年末计数等还可能限制 $t\in\mathbb N$。
\field{预期方法}{确认初始量、时间单位、倍率与时刻起点，再代入或逐次比较阈值。}
\field{常见错误}{把年折旧10\%写成乘0.1；“超过”用成“达到”；时间起点错一位；把衰减至接近0当实际为0。}
\field{达标证据}{模型还原已知数据，结果满足时间范围、整数与精度要求。}
{\small 《教材帮》书内p.131、132、154；PDF第140、141、163页。关联9题；例如：教材p.106 4.2 思考 未编号题组；教材p.108 4.2 例题 例题1；教材p.108 4.2 例题 例题2}\par\end{minipage}\par
\par\medskip\noindent\begin{minipage}{\linewidth}\section{复习顺序与范围边界}
复习顺序建议：先核查集合元素、论域、定义域这些对象条件；再练表示法互换和函数求值；继而用任意两点作差证明单调性，用定义域对称与恒等式判奇偶；最后以配方、基本不等式和指数单调性处理最值、参数与模型。每道题核对“条件—结论—方法—端点”四项。
本次不新加对数及其运算、对数函数、导数、三角函数。教材原题里出现三角形或扇形等初中几何背景，可调用题给公式和初中几何，不据此扩张考查范围。指数模型中的 $\mathrm e$ 可按给定正常数与幂运算处理，不必引入对数。
教师复核时优先查：p.58甲车是否超速的条件不足；p.72插值与平滑读图的非唯一性；整数产量和价格档位；分母禁点；重根与降次；严格不等号；奇偶定义域；离散指数模型的起始时刻。
\end{minipage}
\end{document}
